\chapter{Materiais e métodos}\label{cap:metodos}

\section{Pesquisa bibliográfica e bases de dados}

\TextoModelo{Indique bases consultadas, critérios de seleção, mapas, imagens, modelos digitais de elevação, dados geoquímicos e demais produtos utilizados.}

\section{Etapa de campo}

\TextoModelo{Informe período, escala de trabalho, equipamentos, critérios de descrição de afloramentos, amostragem, controle de qualidade e registro de coordenadas.}

\begin{quadro}
\caption{Exemplo de organização do levantamento de campo}\label{qua:campo}
\centering
\begin{tabularx}{\textwidth}{>{\bfseries}l X l}
\toprule
Atividade & Registro mínimo & Produto \\
\midrule
Descrição de afloramento & Litologia, textura, estruturas, alteração e fotografias & Ficha de ponto \\
Medida estrutural & Tipo de estrutura, atitude, método e qualidade & Banco estrutural \\
Amostragem & Código, material, finalidade e posição & Cadeia de custódia \\
\bottomrule
\end{tabularx}
\FonteTCC{Elaboração própria}
\end{quadro}

\section{Etapa laboratorial}

\TextoModelo{Descreva preparação das amostras, métodos analíticos, equipamentos, padrões, limites de detecção, brancos, duplicatas e tratamento de incertezas.}

\section{Tratamento e integração dos dados}

\TextoModelo{Explique os procedimentos estatísticos, petrográficos, geoquímicos, geocronológicos, geofísicos, geotécnicos ou de geoprocessamento adotados. Declare softwares e versões quando forem relevantes à reprodutibilidade.}
